\section{Cost-directed construction and implementation evidence}\label{sec:study50}
We now execute the distinction between constructing a safe policy and improving its actual economic performance. The new policies use the same nonlinear investment primitives as Section~\ref{sec:study49}; neither the price parameter nor the terminal penalty is changed to favor the proposed action rule. All earlier continuation weights, owners, policies, clocks and adverse contrasts retain their historical identities. The recovered R50 main experiment and common-accuracy amendment had separate local pre-execution freezes. R51 is an explicitly known-design reproduction, with its own Git source identity and runner clocks; the same draws do not provide additional independent observations.

\subsection{Prospective resource selection and a residual-driven comparator}
For the primary block, select powers of two $N\in\{4,8,16,32,64\}$ and $K,M\in\{1,2,4,8,16\}$ before constructing a policy. Among allocations satisfying the primitive witness budget $a/N+b/K+q/M\le\varepsilon-10^{-8}$, minimize the exact integer query proxy $T(N+1)^d(K+1)M$, with lexicographic ties. The reserve is checked against the realized certificate; it is not a universal fixed-precision theorem. This rule supplies an executed version of the separated allocation, rather than a retrospective proxy calculation. Conventional FVI receives the same allocation but its own actual-modulus certificate. The witness planning formula alone does not guarantee FVI attainment.

The catalogue contains seven economic cells: $d=2,3,4$, $T=2,p=1,\varepsilon=5$, and $d=2$, $T=2,3$, $p=1,4,\varepsilon=2$. Witness and uniform FVI are each run on all cells. Two additional services choose the smallest isotropic power-of-two allocation $N=K=M$ satisfying the same primitive target-five budget in $d=2,T=2,p=1$. A further service in that cell uses residual-driven FVI. Each of these seventeen specifications is executed three times in separate single-CPU processes, with identical warm-up and rotated service order. Repetitions estimate local timing dispersion, not a distribution of economic tasks or neural initializations.

The residual-driven comparator queries its own future Bellman objective at interval endpoints and midpoints. It inserts coordinate knots where the largest interpolation surplus across the other coordinates is greatest, then equidistributes square-root surplus mass along each axis, with dyadic rounding and explicit monotone-knot checks. This last step follows the local quadratic interpolation-error scaling, rather than the curvature-mass rule of the earlier graded block. All pilot queries, residual comparisons, compilation and output are charged. Recorded geometry, not the method name, determines whether adaptation occurred. This is a coordinate-tensor comparator, not a claim to have executed the whole adaptive sparse-grid literature.

The original R50 design and source were frozen locally before their measured execution. Its early bisection-only development test and subsequent equidistribution correction are retained with their original records. R51 freezes the recovered algorithm before a known-design reproduction; it is not presented as a new blind experiment. The present clocks belong only to this execution, and the earlier R49 evidence is unchanged.

{\small
\setlength{\tabcolsep}{3pt}
\begin{longtable}{@{}llrrrrrr@{}}
\caption{Primary prospective allocations and realized certificates}\label{tab:services50}\\
\toprule
$d;T;p$ & Method & $N,K,M$ & $\varepsilon$ & Bound & Pass & Sec. & Sec. range \\
\midrule
\endfirsthead
\toprule
$d;T;p$ & Method & $N,K,M$ & $\varepsilon$ & Bound & Pass & Sec. & Sec. range \\
\midrule
\endhead
\bottomrule
\endfoot
2;2;1 & W & 32,2,2 & 2 & 1.9713 & 3/3 & 0.0847 & [0.0841,0.0852] \\
2;2;1 & W I & 8,8,8 & 5 & 4.6296 & 3/3 & 0.0195 & [0.0193,0.0200] \\
2;2;1 & W & 16,1,1 & 5 & 3.9427 & 3/3 & 0.0240 & [0.0238,0.0268] \\
2;2;4 & W & 32,8,1 & 2 & 1.9317 & 3/3 & 0.1086 & [0.1072,0.1105] \\
2;3;1 & W & 64,4,2 & 2 & 1.7567 & 3/3 & 0.5059 & [0.5001,0.5075] \\
2;3;4 & W & 64,4,2 & 2 & 1.9149 & 3/3 & 0.5058 & [0.4962,0.5095] \\
3;2;1 & W & 16,1,1 & 5 & 4.0581 & 3/3 & 0.4424 & [0.4335,0.4610] \\
4;2;1 & W & 8,8,2 & 5 & 4.8628 & 3/3 & 1.8441 & [1.8348,1.9009] \\
2;2;1 & A & 16,1,1 & 5 & 4.2586 & 3/3 & 1.0709 & [1.0636,1.0941] \\
2;2;1 & F & 32,2,2 & 2 & 2.1589 & 0/3 & 0.0644 & [0.0631,0.0646] \\
2;2;1 & F I & 8,8,8 & 5 & 5.3703 & 0/3 & 0.0232 & [0.0227,0.0233] \\
2;2;1 & F & 16,1,1 & 5 & 4.2199 & 3/3 & 0.0243 & [0.0242,0.0246] \\
2;2;4 & F & 32,8,1 & 2 & 2.0923 & 0/3 & 0.0900 & [0.0869,0.0918] \\
2;3;1 & F & 64,4,2 & 2 & 1.7761 & 3/3 & 0.3650 & [0.3495,0.3661] \\
2;3;4 & F & 64,4,2 & 2 & 1.9038 & 3/3 & 0.3606 & [0.3590,0.3723] \\
3;2;1 & F & 16,1,1 & 5 & 4.3316 & 3/3 & 0.3890 & [0.3831,0.4258] \\
4;2;1 & F & 8,8,2 & 5 & 5.653 & 0/3 & 2.4360 & [2.4275,2.4645] \\
\end{longtable}
}
\noindent{\footnotesize W: native witness with equivalent ReLU realization; F: uniform FVI; A: residual-driven FVI; I: isotropic allocation. Each specification has three identical mathematical checkpoints and three distinct clocks. The primitive planning guarantee is witness-specific; FVI uses its actual critic modulus and certificate. Decimal rounding never determines attainment.\par}

The primary block certifies 24 of 24 witness services, 12 of 24 uniform FVI services and 3 of three residual-driven services. These counts include timing repetitions, not independent policy draws. The residual comparator has nonuniform axes in all 3 date models of its repeated checkpoint. Its bound is 4.2586, compared with 4.2199 for uniform FVI at the same quota. Its pilot work is retained.

In the isotropic comparison cell, the planned witness uses 1,156 target innovation evaluations, versus 11,664 for the isotropic service. The latter has fewer state nodes; query savings are therefore not a claim of less work in every category. The table reports measured complete clocks rather than inferring a speed ratio from this count.


\subsection{A separate common-accuracy completion}
Inspection of the primary block showed that sharing a witness-sufficient allocation does not make the FVI certificate pass. We therefore froze a separate, explicit service: in $d=2,3,4$ at $T=2,p=1,\varepsilon=5$, begin with the same planner; if a method's actual certificate fails, double only $N$, reconstruct from primitives and retain the failed attempt. Stop at first success or at $N=64$. Both methods face this rule. There are three new isolated repetitions per method and dimension. The source freeze and results are separate from the primary catalogue, and old failures remain failures.

{\small
\setlength{\tabcolsep}{3pt}
\begin{longtable}{@{}rlrrrrrrr@{}}
\caption{Separate common-target completion: horizon two, price one, target five}\label{tab:common50}\\
\toprule
$d$ & Method & $N$ range & Tries & Bound & Pass & Sec. & Sec. range & MiB \\
\midrule
\endfirsthead
\toprule
$d$ & Method & $N$ range & Tries & Bound & Pass & Sec. & Sec. range & MiB \\
\midrule
\endhead
\bottomrule
\endfoot
2 & W & 16--16 & 1 & 3.9427 & 3/3 & 0.0242 & [0.0223,0.0281] & 41.4 \\
2 & F & 16--16 & 1 & 4.2199 & 3/3 & 0.0244 & [0.0236,0.0275] & 41.7 \\
3 & W & 16--16 & 1 & 4.0581 & 3/3 & 0.4383 & [0.4155,0.4386] & 44.3 \\
3 & F & 16--16 & 1 & 4.3316 & 3/3 & 0.3913 & [0.3840,0.3935] & 42.3 \\
4 & W & 8--8 & 1 & 4.8628 & 3/3 & 1.7418 & [1.7139,1.8030] & 50.5 \\
4 & F & 8--16 & 2 & 3.1181 & 3/3 & 29.0946 & [28.9063,29.3447] & 169.8 \\
\end{longtable}
}
\noindent{\footnotesize All failed state refinements are included in the complete prefix. MiB is the maximum process peak over three repetitions, not just model storage. The amendment followed the primary outcomes and was frozen before its own execution. Neither block overwrites the other. Host exclusivity and CPU frequency were not controlled.\par}

The returned policies receive new direct actual-cost comparisons below. These observations provide a common target in every reported dimension. They remain a finite low-dimensional tensor experiment. The construction and corner factors $(N+1)^d$ and $2^d$ are not removed by the compiler, the terminal repair or this allocation rule.

\subsection{Direct policy costs and the source of removable loss}
The primary cost experiment fixes twenty-five groups. Sixteen use the original first-target-two R49 witness and FVI policies in all four price--horizon cells, evaluated under the uniform law and each of the three declared point laws. Seven use the new planned services under uniform initial states; one compares residual-driven and uniform FVI; one uses the isotropic allocations. A separate three-group family evaluates the common-accuracy amendment's returned policies. Every group contains both original and final-repaired policies, with the same eight-bit observation and $1/4096$ action contract. Both generators receive the same improvement.

Each group uses 131,072 continuous-law interval paths and analytic final-innovation integration. Four absolute costs and all six pair contrasts are recorded. The primary 250 estimands and supplemental 30 estimands are covered by the conservative endpoint union-bound budget of 512 and log upper bound 13. The supplemental groups have their own conditional independent-bin model within the declared simultaneous family; the policies may depend on earlier construction information. Fixed streams establish reproducibility only. The older R49 family retains its own separate error account.

{\small
\setlength{\tabcolsep}{3pt}
\begin{longtable}{@{}lrrr@{}}
\caption{Actual policy differences and removable final-action loss}\label{tab:policy50}\\
\toprule
Case & $J_L-J_R$ & $J_{L^+}-J_{R^+}$ & $J_L-J_{L^+}$ \\
\midrule
\endfirsthead
\toprule
Case & $J_L-J_R$ & $J_{L^+}-J_{R^+}$ & $J_L-J_{L^+}$ \\
\midrule
\endhead
\bottomrule
\endfoot
C2;2;1;U & [0.023033,0.028988] & [0.012051,0.017775] & [0.011448,0.014567] \\
C3;2;1;U & [0.01796,0.023776] & [0.010129,0.015794] & [0.0077535,0.010757] \\
C4;2;1;U & [0.033636,0.039473] & [0.015912,0.021433] & [0.016432,0.019542] \\
I2;2;1;U & [0.042713,0.048868] & [0.018209,0.023856] & [0.023541,0.026849] \\
O2;2;1;M & [0.0064146,0.011415] & [-0.0023843,0.0023843] & [0.007622,0.010239] \\
O2;2;1;L & [0.054972,0.060123] & [0.044007,0.048792] & [0.0099411,0.012722] \\
O2;2;1;H & [0.0042173,0.0090667] & [0.0022567,0.007049] & [0.0014934,0.003946] \\
O2;2;1;U & [0.013161,0.018416] & [0.0064562,0.011605] & [0.0056474,0.0083165] \\
O2;2;4;M & [0.035287,0.040496] & [0.030381,0.035526] & [0.0036933,0.0063836] \\
O2;2;4;L & [0.020345,0.025506] & [0.017976,0.023095] & [0.0010828,0.0036974] \\
O2;2;4;H & [0.039236,0.044497] & [0.034549,0.039692] & [0.0034416,0.006155] \\
O2;2;4;U & [0.011181,0.016658] & [0.0058295,0.011256] & [0.004188,0.0069254] \\
O2;3;1;M & [0.015697,0.021868] & [0.012485,0.018641] & [0.0016956,0.0047777] \\
O2;3;1;L & [0.026367,0.032516] & [0.022928,0.029055] & [0.0020031,0.005116] \\
O2;3;1;H & [0.011221,0.017268] & [0.0095515,0.015586] & [0.00030976,0.0033545] \\
O2;3;1;U & [0.0095127,0.015773] & [0.0065302,0.012761] & [0.0014967,0.0045926] \\
O2;3;4;M & [0.010433,0.016948] & [0.0089337,0.015429] & [9.4671e-05,0.0033779] \\
O2;3;4;L & [0.0095446,0.01608] & [0.0074164,0.013916] & [0.0005048,0.0037993] \\
O2;3;4;H & [0.0073415,0.013898] & [0.0046083,0.011119] & [0.0011154,0.0044219] \\
O2;3;4;U & [0.0056055,0.012273] & [0.0039105,0.010565] & [0.00021625,0.0035063] \\
Q2;2;1;U & [0.01186,0.017187] & [0.0056983,0.010907] & [0.0055893,0.0083122] \\
P2;2;1;U & [0.022852,0.028801] & [0.011835,0.01755] & [0.011497,0.014615] \\
Q2;2;4;U & [0.0086864,0.014074] & [0.004022,0.0093656] & [0.003535,0.0062382] \\
Q2;3;1;U & [0.01136,0.017753] & [0.0081953,0.014553] & [0.0017403,0.0048791] \\
Q2;3;4;U & [0.010559,0.017349] & [0.0074835,0.014248] & [0.0020091,0.0053568] \\
P3;2;1;U & [0.017874,0.023687] & [0.0099988,0.015657] & [0.0078066,0.010814] \\
P4;2;1;U & [0.033196,0.039029] & [0.015676,0.021193] & [0.016374,0.019482] \\
A2;2;1;U & [-0.0024517,0.0023493] & [-0.0024125,0.0023985] & [0.00062241,0.0030972] \\
\end{longtable}
}
\noindent{\footnotesize A case is block/dimension; horizon; price; initial law. O: original R49 target-two policies; P: primary target-five plan; Q: primary target-two plan; C: common-accuracy amendment; I: isotropic; A: surplus FVI versus uniform FVI. Except A, left is witness and right is uniform FVI. U is uniform; L,M,H are point states 1/8,1/2,7/8. Superscript + denotes the identical terminal safety rule. Five-significant-digit display is not an exact endpoint; the rational records govern all conclusions.\par}

All 27 declared witness comparison groups identify a strictly positive reduction of the witness incumbent's actual expected cost. The simultaneous gain lower endpoints range from 9.4671e-05 to 0.023541; upper endpoints range from 0.0033545 to 0.026849. This is a law-specific statement for the stated policies, not strict improvement at every initial state or 27 independent training successes.

Before repair, witness cost is identified as higher in 27 groups. After both methods receive the same repair, it remains higher in 26, is lower in 0, and is unresolved in 1. The residual-driven FVI comparison is unresolved both before and after repair. Thus the terminal construction repairs a genuine economic loss without manufacturing a reversal of the strong conventional comparison.

The active evaluator records 3,670,016 interval paths, 7,340,032 cell-gate queries and 5,447,560 accepted changes across the two incumbents. No acquisition ambiguity or broad-hull fallback is recorded. Its joint direct-evaluation time before durable record output is 127.95 seconds. That total is evaluation work, not a speed comparison between independently executed policies.


For each incumbent the direct reduction is the centered final contrast. The remaining witness--FVI difference after both repairs is an actual comparison of the two changed policies. The identity
\[
 J_W-J_F=(J_W-J_{W^+})+(J_{W^+}-J_{F^+})-(J_F-J_{F^+})
\]
separates removable final-action losses from the difference that remains. It does not attribute every remaining discrepancy to a single representation mechanism. The deposited date-level action means, terminal state means, accepted and blocked proposals, and capacity/quantum repair counts permit this distinction to be checked. A finite-sample positive gain can justify a strictly positive replacement charge up to its simultaneous lower endpoint; a deterministic nonincrease alone does not justify an arbitrary implementation fee.

\subsection{Actual-cost frontiers and complete resource boundaries}
The supplement reports every group's four absolute cost intervals and every direct contrast. For each new construction, its resource coordinate includes the full own-future service and, separately, the joint evaluation charge. The reader can use either a shared fixed evaluation cost for the catalogue decision or the complete joint cost required to certify a particular comparison. We do not claim to have measured independent marginal costs for four separately simulated policies when the implementation shares their earlier trajectories. The pair clock is an explicitly charged object, not an unpriced convenience.

{\small
\setlength{\tabcolsep}{3pt}
\begin{longtable}{@{}llrrrr@{}}
\caption{Construction work and actual cost of new implementations}\label{tab:joint50}\\
\toprule
Case & Method & Build sec. & Original cost & Repaired cost & Pair sec. \\
\midrule
\endfirsthead
\toprule
Case & Method & Build sec. & Original cost & Repaired cost & Pair sec. \\
\midrule
\endhead
\bottomrule
\endfoot
C2;2;1;U & W & 0.0242 & [0.6242,0.6378] & [0.61117,0.62482] & 3.799 \\
C2;2;1;U & F & 0.0244 & [0.59821,0.61177] & [0.59628,0.60988] & 3.799 \\
C3;2;1;U & W & 0.4383 & [0.58992,0.60068] & [0.58062,0.59146] & 6.102 \\
C3;2;1;U & F & 0.3913 & [0.56904,0.57982] & [0.56768,0.57849] & 6.102 \\
C4;2;1;U & W & 1.7418 & [0.62188,0.63199] & [0.60399,0.61391] & 9.790 \\
C4;2;1;U & F & 29.0946 & [0.5855,0.59525] & [0.5854,0.59515] & 9.790 \\
I2;2;1;U & W & 0.0195 & [0.63959,0.65386] & [0.61456,0.6285] & 3.744 \\
I2;2;1;U & F & 0.0232 & [0.59407,0.60781] & [0.59363,0.60737] & 3.744 \\
Q2;2;1;U & W & 0.0847 & [0.60972,0.62332] & [0.60276,0.61638] & 3.870 \\
Q2;2;1;U & F & 0.0644 & [0.59518,0.60882] & [0.59444,0.6081] & 3.870 \\
P2;2;1;U & W & 0.0240 & [0.6222,0.63577] & [0.60912,0.62274] & 3.901 \\
P2;2;1;U & F & 0.0243 & [0.59639,0.60993] & [0.59445,0.60803] & 3.901 \\
Q2;2;4;U & W & 0.1086 & [0.69243,0.70702] & [0.68754,0.70213] & 3.919 \\
Q2;2;4;U & F & 0.0900 & [0.68106,0.69562] & [0.68086,0.69542] & 3.919 \\
Q2;3;1;U & W & 0.5059 & [0.72237,0.73696] & [0.71906,0.73365] & 4.406 \\
Q2;3;1;U & F & 0.3650 & [0.70782,0.7224] & [0.70769,0.72228] & 4.406 \\
Q2;3;4;U & W & 0.5058 & [0.86908,0.88468] & [0.86539,0.881] & 4.513 \\
Q2;3;4;U & F & 0.3606 & [0.85511,0.87074] & [0.85451,0.87015] & 4.513 \\
P3;2;1;U & W & 0.4424 & [0.58981,0.60055] & [0.58046,0.59128] & 6.294 \\
P3;2;1;U & F & 0.3890 & [0.56902,0.57978] & [0.56764,0.57844] & 6.294 \\
P4;2;1;U & W & 1.8441 & [0.61951,0.62957] & [0.60168,0.61155] & 7.713 \\
P4;2;1;U & F & 2.4360 & [0.58357,0.59328] & [0.58332,0.59303] & 7.713 \\
A2;2;1;U & A & 1.0709 & [0.59772,0.61129] & [0.59585,0.60945] & 3.797 \\
A2;2;1;U & F & 0.0243 & [0.59777,0.61134] & [0.59585,0.60946] & 3.797 \\
\end{longtable}
}
\noindent{\footnotesize Pair sec. is a shared cost of the whole four-policy comparison, repeated in the display for clarity, not independently attributable to each row. Construction medians use the same cell and fresh records; no historical clocks are spliced. Absolute-cost intervals are wider than centered pair intervals. Charges enter Proposition~\ref{prop:joint50} with their explicitly chosen shared or marginal accounting convention.\par}

All stored coefficients, original owners, action scalars, model evaluation categories, transform relaxations, interpolation/corner operations, terminal derivative queries, whole-cell tests, exact fallbacks, numerical interval widths, serialized bytes and process peak resident memory remain in the records. The counters are primitive categories, not a complete FLOP or bit-operation count. The latter would require an additional cost model. Direct comparison time is not added to clocks from another host or to a historical failure in order to manufacture a new complete service.

The retained sensor catalogue has a structural interpretation through equation~\eqref{eq:technology50}. Its prices remain normalized primitives, and its confidence-envelope regret is relative to the best observed implementation in the finite catalogue. This revision neither estimates a sensor technology nor substitutes a sample-mean winner for the true expected-net-cost minimizer.

\subsection{Historical findings that the new evidence does not replace}
The earlier 93 of 96 strictly higher witness-cost contrasts remain visible in Table~\ref{tab:direct49}. Original curvature-FVI never changed its grid in the realized block. The later genuinely graded method also remains an adverse conventional-comparator experiment: it was less precise in all twelve matched final services and slower in all twelve. These two findings concern the two specified heuristics, not all adaptive dynamic programming. The present surplus comparator adds an independently recorded realized partition, while uniform FVI remains the strong baseline.

Likewise, a successful final-action repair does not establish better earlier continuation learning, universal neural superiority, a guarantee for unrestricted nonconvex training, or a dimension-free policy construction. What is established is a cost-directed feasible improvement with an explicit safety condition, its executed original-economy specialization, and source-bound resource and actual-cost comparisons. These results extend the NBO policy-evaluation and improvement program without replacing its economic subject or erasing its previous evidence.
